\section{Dynamic tests and requirement-grounded witnesses}\label{sec:C25:testjack}
TestJack audits benchmark-passing coding trials by generating tests tied to prompt requirements, checking them against a reference solution, and reviewing whether failures represent genuine requirement violations~\cite{yao2026testjack}. Its distinction between insufficient tests and underspecified prompts is essential: disagreement with the reference is not enough to refute a candidate. The paper reports substantial refutations in its audited passing-trial corpus; that proportion is not a prevalence estimate for all coding-agent executions.

The method's witness structure has three obligations. The test must correspond to a required behavior, pass on the reference, and fail on the candidate. The executable conditions are reproducible, but the requirement correspondence remains interpretive. A faulty reference or ambiguous prompt can invalidate the inference. A reviewer agent reduces some mistakes but does not become an oracle merely because its role is called reviewer.

A stronger causal ablation keeps the candidate set and execution environment fixed while varying requirement grounding, reference validation, and review. Measure how many proposed witnesses survive independent adjudication and what defect classes they expose. Also examine false refutations: a method that discovers many failures by imposing extra requirements can reduce a leaderboard score while making the evaluation less valid.

Witness sharing across trials can amortize cost. If several candidates fail the same requirement, one validated test can inspect them all. This changes the unit of evidence: outcomes share a task and a discovered witness, so they are statistically dependent. Uncertainty should reflect that clustering. Randomly auditing trials can estimate properties of the declared audit population, but adaptive witness reuse and selective follow-up require a carefully described sampling scheme.

The implementation path should preserve the original repository, candidate patch, reference patch, generated witness, native test runner, and adjudication record. An artifact link or algorithm description does not establish that all pieces are available at an immutable revision. Until the exact implementation is inspected and run, report the method and author results separately from reproduction. This is particularly important for evaluator evolution, where changing a prompt or test-generation model can alter the measured defect population.

